\scanpage{16}
profinis, dans une certaine mesure !) utiliser le langage géométrique des groupoïdes fondamentaux de $\mathbb P^1(\mathbb C)-\{0,1,\infty\}$. On \deleted{pourra} notera donc $\pi=\pi_{0,3}$ le groupe discret
\[
\pi=\pi_{\unsure{1|*}}(\overbrace{\mathbb P^1(\mathbb C)\setminus\{0,1,\infty\}}^{U=U_{0,3}},P)
\]
où $\mathbb P^1(\mathbb C)$ est muni de la structure \inserted{de sphère} riemannienne \inserted{orientée} unique pour laquelle la rotation \inserted{$\dot\rho$} d'angle $+\frac{2\pi}{3}$ autour de l'axe $P,\bar P$ \deleted{\unsure{qui}} induit sur $(0,1,\infty)$ la permutation circulaire $0\longrightarrow1\longrightarrow\infty\longrightarrow0$ (ceci correspond à l'orientation habituelle de la sphère), et où les aut complexes \inserted{$\dot\sigma_0,\dot\sigma_1,\dot\sigma_\infty$} de $\mathbb P^1_{\mathbb C}$ qui \deleted{\unsure{échangent}} \inserted{\deleted{fixent}} induisent \deleted{\unsure{sur}} sur $(0,1,\infty)$ les transpositions fixant resp. $0,1,\infty$, sont \inserted{resp.} les \deleted{rotations} \inserted{symétries} autour des axes \deleted{$\unsure{P,\bar P}$} $(0,2)$, $(1,-1)$ et $(\infty,\frac12)$.
\margin{Sur la sphère de Riemann on a
\[
\begin{aligned}
\sigma_0(z)&=\frac{-z}{1-z}=\frac{z}{z-1}\\
\sigma_1(z)&=\frac1z\\
\sigma_\infty(z)&=1-z\\
\rho(z)&=\frac1{1-z}\\
\rho^2(z)&=1-\frac1z.
\end{aligned}
\]}
Quant à la conjugaison complexe, c'est la symétrie par un plan $(0,1,\infty)$. On considère $\unsure{\mathfrak S_3|\mathcal S_3}$ comme opérant sur $(0,1,\infty)$ dans cet ordre, et $\unsure{\mathfrak S_3|\mathcal S_3}\times\mathbb Z/2$ comme opérant sur la sphère, via les $\dot\sigma_0,\dot\sigma_1,\dot\sigma_\infty\in\unsure{\mathfrak S_3|\mathcal S_3}$ et $\dot\tau\in\mathbb Z/2$. On considère l'extension \inserted{$\unsure{\mathcal E_0}$} de $\unsure{\mathfrak S_3|\mathcal S_3}\times\mathbb Z/2$ par $\pi$, notée \inserted{$\mathcal E_{0,3}$ ou} $\mathcal E$, définie par
\[
\mathcal E=\pi_1\left(\underset{\substack{\Vert\\U_{0,3}}}{U_{\deleted{\unsure{0,3}}}},\underbrace{\unsure{\mathfrak S_3|\mathcal S_3}\times\mathbb Z/2}_{\Gamma};P\right)
\]
\mspage{15}
